\documentclass[12pt]{article}
\usepackage[T1]{fontenc}
\usepackage[utf8]{inputenc}
\usepackage{lmodern}
\usepackage{textcomp}
\usepackage{enumitem}
\usepackage{geometry}
\usepackage{graphicx}
\usepackage{amsmath, amssymb}
\geometry{margin=1in}
\begin{document}
\begin{center}
Economics - Undergraduate Level Assessment 23 \
Total Marks: 40 \
Answer all questions.
\end{center}

\section*{Multiple Choice (10 marks)}
\begin{enumerate}[label=\arabic*.]
\item Which of the following statements is most consistent with a capitalist market economy?
\begin{enumerate}[label=(\alph*)]
    \item Economic resources are allocated according to the decisions of the central bank.
    \item Private property is fundamental to innovation, growth, and trade.
    \item A central government plans the production and distribution of goods.
    \item Most wages and prices are legally controlled.
\end{enumerate}
\item A dark room technician who develops film into photographs loses his job because few people use film cameras any more. This is an example of which of the following?
\begin{enumerate}[label=(\alph*)]
    \item Seasonal unemployment
    \item Casual unemployment
    \item Cyclical unemployment
    \item Structural unemployment
    \item Underemployment
\end{enumerate}
\item If a Johansen "trace" test for a null hypothesis of 2 cointegrating vectors is applied to a system containing 4 variables is conducted, which eigenvalues would be used in the test?
\begin{enumerate}[label=(\alph*)]
    \item The largest and smallest
    \item The smallest and second smallest
    \item The middle two
    \item The smallest 2
    \item The largest 2
\end{enumerate}
\item If real GDP per capita was $10000 in 1990 and $15000 in 2000 then the amount of economic growth is
\begin{enumerate}[label=(\alph*)]
    \item 75 percent.
    \item 0.5 percent.
    \item 20 percent.
    \item 100 percent.
    \item 25 percent.
\end{enumerate}
\item Which of the following would shift the aggregate demand curve to the left?
\begin{enumerate}[label=(\alph*)]
    \item An increase in consumer confidence.
    \item Increase in the prices of goods and services.
    \item Business firms expect lower sales in the future.
    \item A technological advancement in the production process.
    \item An increase in the wealth of consumers.
\end{enumerate}
\end{enumerate}

\section*{True / False (10 marks)}
\textit{Answer True or False}
\begin{enumerate}[label=\arabic*.]
\setcounter{enumi}{5}
\item True or False: Theo's job loss due to the seasonal closure of the pool exemplifies natural unemployment, which includes seasonal job fluctuations.
\item True or False: In a perfectly competitive market, firms must raise wages to attract additional workers due to increased demand for labor.
\item True or False: The velocity of money in the United States, given a total transaction value of $100,000 and a money supply of $10,000, is 20.
\item True or False: VARs often produce better forecasts than simultaneous equation structural models due to their flexibility in capturing dynamic relationships.
\item True or False: Demand for a good with no close substitutes is likely to have a demand curve that is the least elastic.
\end{enumerate}

\section*{Long Form Response (20 marks)}
\textit{Answer each question with a clear and complete explanation.}
\begin{enumerate}[label=\arabic*.]
\setcounter{enumi}{10}
\item Discuss the impact of multicollinearity on the properties of the Ordinary Least Squares (OLS) estimator, particularly regarding its unbiasedness and efficiency.
\item Discuss the characteristics of a kinked demand curve and critically analyze why the relationship between price and marginal cost is significant in this model.
\end{enumerate}

\vfill
\noindent\textit{End of Paper}
\end{document}